\documentclass[../main]{subfiles}
\begin{document}
\chapter{Servomotor}
\img{images/servo/partservo}{Partes del servomotor}
Un servomotor es un dispositivo actuador que tiene la capacidad de ubicarse en cualquier posici\'on{} dentro de su rango de operaci\'on{}, y de mantenerse estable en dicha posici\'on{}. Est\'a{} formado por un motor de corriente continua, una caja reductora y un circuito de control, y su margen de funcionamiento suele ser de menos de una vuelta completa.

\img{images/servo/servomotor}{Servomotor de aeromodelismo.}
Los servos de modelismo se utilizan frecuentemente en sistemas de radiocontrol y en rob\'otica{}, pero su uso no est\'a{} limitado a estos.

El componente principal de un servo es un motor de corriente continua, que realiza la funci\'on{} de actuador en el dispositivo: al aplicarse un voltaje entre sus dos terminales, el motor gira en un sentido a alta velocidad, pero produciendo un bajo par. Para aumentar el par del dispositivo, se emplea una caja reductora, que transforma gran parte de la velocidad de giro en torsi\'on{}.

\video{https://www.youtube.com/watch?v=mk9UkQCeENc}{Funcionamiento de Servomotor.}

\section{Control}
\subsection{Control de posicionado}

\img{images/servo/servo}{Diagrama del circuito de control implementado en un servo.}

El dispositivo utiliza un circuito de control para realizar la ubicaci\'on{} del motor en un punto, consistente en un controlador proporcional.

El punto de referencia o \emph{setpoint}: que es el valor de posici\'on{} deseada para el motor, se indica mediante una se\~nal{} de control cuadrada. El ancho de pulso de la se\~nal{} indica el \'angulo{} de posici\'on{}: una se\~nal{} con pulsos m\'as{} anchos (es decir, de mayor duraci\'on{}) ubicar\'a{} al motor en un \'angulo{} mayor, y viceversa. (PWM)

Inicialmente, un amplificador de error calcula el valor del error de posici\'on{}, que es la diferencia entre la referencia y la posici\'on{} en que se encuentra el motor. Un error de posici\'on{} mayor significa que hay una diferencia mayor entre el valor deseado y el existente, de modo que el motor deber\'a{} rotar m\'as{} r\'apido{} para alcanzarlo; uno menor, significa que la posici\'on{} del motor est\'a{} cerca de la deseada por el usuario, as\'i{} que el motor tendr\'a{} que rotar m\'as{} lentamente. Si el servo se encuentra en la posici\'on{} deseada, el error ser\'a{} cero, y no habr\'a{} movimiento.

\img{images/servo/pwm}{Ejemplos de se\~nales{} de control utilizadas.}

Para que el amplificador de error pueda calcular el error de posici\'on{}, debe restar dos valores de voltaje anal\'ogicos. La se\~nal{} de control PWM se convierte entonces en un valor anal\'ogico{} de voltaje, mediante un convertidor de ancho de pulso a voltaje. El valor de la posici\'on{} del motor se obtiene usando un potenci\'ometro{} de realimentaci\'on{} acoplado mec\'anicamente{} a la caja reductora del eje del motor: cuando el motor rote, el potenci\'ometro{} tambi\'e{}n lo har\'a{}, variando el voltaje que se introduce al amplificador de error.

Una vez que se ha obtenido el error de posici\'on{}, este se amplifica con una ganancia, y posteriormente se aplica a los terminales del motor.

\subsection{Utilizaci\'on{}}
Dependiendo del modelo del servo, la tensi\'on{} de alimentaci\'on{} puede estar comprendida entre los \qty{4}{V} y \qty{8}{V}. El control de un servo se reduce a indicar su posici\'on{} mediante una se\~nal{} cuadrada de voltaje: el \'angulo{} de ubicaci\'on{} del motor depende de la duraci\'on{} del nivel alto de la se\~nal{}.

Dependiendo de la marca y modelo utilizado, cada servo tiene sus propios m\'argenes{} de operaci\'on{}. Por ejemplo, para algunos servos los valores de tiempo de la se\~nal{} en alto est\'an{} entre \qty{1}{ms} y \qty{2}{ms}, que posicionan al motor en ambos extremos de giro ($0^\circ$ y $180^\circ$, respectivamente). Los valores de tiempo de alto para ubicar el motor en otras posiciones se hallan mediante una relaci\'on{} completamente lineal: el valor \qty{1,5}{ms} indica la posici\'on{} central, y otros valores de duraci\'on{} del pulso dejar\'ian{} al motor en la posici\'on{} proporcional a dicha duraci\'on{}.

Es sencillo notar que, para el caso del motor anteriormente mencionado, la duraci\'on{} del pulso alto para conseguir un \'angulo{} de posici\'on{} $\phi$ estar\'a{} dado por la f\'ormula{}:
\begin{equation}
t = 1 + \frac{\phi}{180}
\end{equation}

D\'onde{} $t$ est\'a{} dado en milisegundos y $\phi$ en grados. Sin embargo, debe tenerse en cuenta que ning\'un{} valor de \'angulo{} o de duraci\'on{} de pulso puede estar fuera del rango de operaci\'on{} del dispositivo: en efecto, el servo tiene un l\'imite{} de giro de modo que no puede girar m\'as{} de cierto \'angulo{} en un mismo sentido debido a la limitaci\'on{} f\'isica{} que impone el potenci\'ometro{} del control de posici\'on{}.

Para bloquear el servomotor en una posici\'on{}, es necesario enviarle continuamente la se\~nal{} con la posici\'on{} deseada. De esta forma, el sistema de control seguir\'a{} operando, y el servo conservar\'a{} su posici\'on{} y se resistir\'a{} a fuerzas externas que intenten cambiarlo de posici\'on{}. Si los pulsos no se env\'ian{}, el servomotor quedar\'a{} liberado, y cualquier fuerza externa puede cambiarlo de posici\'on{} f\'acilmente.

\actividad{Contestar: en base al texto}
\begin{enumerate}
\item \textquestiondown{}Qu\'e{} es un servomotor?
\item \textquestiondown{}Cu\'ales{} son los componentes principales de un servomotor?
\item \textquestiondown{}Cu\'al{} es la funci\'on{} del motor de corriente continua en un servomotor?
\item \textquestiondown{}Qu\'e{} es una caja reductora y qu\'e{} papel desempe\~na{} en un servomotor?
\item \textquestiondown{}C\'omo{} se controla la posici\'on{} de un servomotor?
\item \textquestiondown{}Qu\'e{} indica el ancho de pulso de la se\~nal{} de control en un servomotor?
\item \textquestiondown{}C\'omo{} se calcula el error de posici\'on{} en un servomotor?
\item \textquestiondown{}C\'omo{} se convierte la se\~nal{} de control PWM en un valor anal\'ogico{} de voltaje?
\item \textquestiondown{}C\'omo{} se utiliza un potenci\'ometro{} en un servomotor?
\item \textquestiondown{}Qu\'e{} sucede si no se env\'ian{} pulsos de control a un servomotor?
\end{enumerate}
\end{document}